\documentclass[11pt,a4paper]{article}
\usepackage[utf8]{inputenc}
\usepackage[T1]{fontenc}
\usepackage[french]{babel}
\usepackage{lmodern}
\usepackage[margin=2.5cm]{geometry}
\usepackage{amsmath,amssymb}
\usepackage{booktabs}
\usepackage[hidelinks]{hyperref}

\newcommand{\code}[1]{\texttt{#1}}

\title{Projet MLSF -- Un solveur SAT pour Meowdoku\\ \large Partie II -- Un langage de contraintes compilé vers SAT}
\author{Pablo Bosch \and Cécile Chevrière}
\date{Novembre 2026}

\begin{document}
\maketitle

\paragraph{Organisation du rendu.} Le compilateur se trouve dans le dossier \code{partie2/}.
Il réutilise tel quel le solveur de la partie I (\code{partie1/Session de travail/}). Chaque
étape de la chaîne correspond à un module :
\begin{center}
\begin{tabular}{lll}
\toprule
étape & module & question \\
\midrule
texte $\to$ arbre de syntaxe abstraite (AST) & \code{syntax.py} & Q2 \\
table des symboles, vérifications & \code{symboles.py}, \code{analyse.py} & Q3 \\
déroulement, simplifications $\to$ contraintes sans \code{forall} & \code{deroulement.py}, \code{ir.py} & Q3 \\
encodage des domaines, comparaisons, globales, Tseitin $\to$ FNC & \code{fnc.py} & Q1, Q4 \\
DIMACS, solveur de la partie I, décodage & \code{solveur.py} & Q1, Q5 \\
ligne de commande & \code{csp.py} & \\
\bottomrule
\end{tabular}
\end{center}
L'interface est celle conseillée par l'énoncé (depuis \code{partie2/}) :
\code{python csp.py m.csp} écrit la FNC au format DIMACS sur la sortie standard,
\code{--dimacs f.cnf} l'écrit dans un fichier, et \code{--solve} affiche \code{s SATISFIABLE}
suivi d'une ligne \code{nom = valeur} par variable, ou \code{s UNSATISFIABLE}. Les tests se
lancent par \code{python -m tests.test\_X} pour X parmi \code{syntax}, \code{analyse},
\code{fnc}, \code{exemples}, \code{meowdoku}. Le dossier \code{q6\_meowdoku/} contient les
programmes de la question 6.

%======================================================================
\section{Question 1 -- Du fragment propositionnel au solveur}

La première chaîne complète a été construite pour le fragment propositionnel, puis
généralisée au fil des questions : c'est la même chaîne qui traite aujourd'hui tout le
langage. Le fragment y est un cas particulier, celui où toutes les variables sont booléennes.

\paragraph{Table des symboles.} \code{constraint w;} est rejeté si \code{w} n'est pas déclaré
(\og nom non déclaré 'w' \fg). La table a été étendue à la question 3.

\paragraph{Mise en FNC.} \code{a -> b} est réécrit en $\neg a \vee b$, et la négation est
poussée jusqu'aux littéraux par les lois de De Morgan. On applique ensuite la transformation
de Tseitin.
\begin{itemize}
  \item Une conjonction au sommet donne une clause par conjoint. Une disjonction de littéraux
        est déjà une clause : \code{a -> b}, \code{not (a and b)}, \code{(a and b) -> c} ne
        coûtent aucune variable auxiliaire.
  \item Toute autre sous-formule $\varphi$ reçoit une variable auxiliaire $a$, avec les
        clauses de $a \leftrightarrow \varphi$. Pour $\varphi = l_1 \wedge \dots \wedge l_n$ :
        $\neg a \vee l_i$ pour chaque $i$, et $a \vee \neg l_1 \vee \dots \vee \neg l_n$. Pour
        une disjonction, les clauses symétriques.
  \item Les clauses en double ou tautologiques sont supprimées.
\end{itemize}

\paragraph{Noms des variables.} Le fichier DIMACS commence par des commentaires
\code{c var 3 z} pour les booléens, ou \code{c var 12 grid[1][3] = 3} pour les domaines finis
(question 4). Les variables auxiliaires n'ont pas de nom et ne sont pas affichées.

\paragraph{Validation (\code{tests/test\_fnc.py}).} Un générateur produit des modèles
aléatoires (3 booléens, 1 à 3 contraintes de profondeur 1 à 3, avec \code{not}, \code{and},
\code{or}, \code{->}, \code{=} et \code{!=} entre booléens et les contraintes globales). Un
évaluateur de l'AST, écrit dans le test indépendamment du compilateur, fournit la réponse par
recherche exhaustive. Deux propriétés sont vérifiées :
\begin{enumerate}
  \item la FNC est satisfiable si et seulement si le modèle l'est ;
  \item la restriction d'un modèle de la FNC aux variables nommées est une solution du modèle.
\end{enumerate}
300 modèles sont testés. Le test étendu aux domaines finis (question 4) en ajoute 300 autres.

%======================================================================
\section{Question 2 -- Analyse syntaxique}

\paragraph{Grammaire.} La grammaire (\code{syntax.py}) reste LALR(1)
(\code{parser="lalr"}, analyseur lexical \code{basic}). Les priorités sont encodées par
niveaux successifs, du moins prioritaire au plus prioritaire :
\begin{center}
\code{->} (droite) $<$ \code{or} (gauche) $<$ \code{and} (gauche) $<$ \code{not} $<$
comparaisons (non associatives) $<$ \code{+ -} (gauche) $<$ \code{*} (gauche) $<$
\code{-} unaire $<$ atomes.
\end{center}
\begin{itemize}
  \item Les comparaisons sont plus prioritaires que \code{not} : \code{not x = 2} se lit
        \code{not (x = 2)}. Elles ne s'enchaînent pas, car la règle est
        \code{comparison: sum OP sum}, sans récursion : \code{a = b = c} est une erreur de
        syntaxe.
  \item Les atomes sont un entier, un nom, un accès \code{t[e]...[e]}, un appel
        \code{f(e, ..., e)} ou une expression parenthésée. Un indice est une expression ou le
        joker \code{*}. Aucune expression ne commence par \code{*} : après \code{[}, le
        lexème \code{*} ne peut être que le joker, d'où l'absence de conflit avec la
        multiplication.
  \item Les contraintes globales et \code{abs} sont de simples appels de fonction : en
        ajouter une ne change pas la grammaire.
  \item Un domaine est \code{bool}, \code{sum .. sum} ou un nom (type énuméré). Le conflit
        potentiel entre \og nom de type \fg{} et \og nom utilisé comme borne \fg{}
        (\code{var c : Color;} contre \code{var x : N..M;}) est tranché par le lexème
        suivant : \code{;} dans le premier cas, \code{..} ou un opérateur dans le second.
        Lark n'a signalé aucun conflit.
  \item Une dimension de tableau est une expression ; un nom de type énuméré
        (\code{used[Color]}) est reconnu à l'analyse.
  \item Le corps d'un \code{forall} ne contient que des contraintes et des \code{forall} :
        une déclaration dans une boucle est une erreur de syntaxe.
\end{itemize}
Les onze mots-clés sont réservés grâce à l'analyseur lexical \code{basic}, qui produit
toujours un lexème mot-clé et jamais un \code{NAME} (\code{var in : bool;} est rejeté). Un nom
qui commence par un mot-clé (\code{index}, \code{order}) reste un nom.

\paragraph{Arbre de syntaxe abstraite.} Une classe par construction : \code{ConstDecl},
\code{ConstArrayDecl}, \code{EnumDecl}, \code{VarDecl}, \code{ArrayDecl}, \code{Constraint},
\code{Forall}, \code{RangeDomain}, \code{EnumDomain} (\code{bool} compris), et pour les
expressions \code{Var}, \code{IntLit}, \code{Not}, \code{And}, \code{Or}, \code{Implies},
\code{Compare}, \code{BinOp}, \code{Neg}, \code{Access}, \code{Wildcard}, \code{Call}. Les
classes du code de départ sont conservées. \code{pretty} affiche l'arbre en
parenthésant chaque opérateur ; \code{source} réécrit le modèle en texte entièrement
parenthésé.

\paragraph{Tests (\code{tests/test\_syntax.py}).} Les tests couvrent les priorités et
associativités (anciens et nouveaux opérateurs), les déclarations, 13 erreurs de syntaxe
(avec leur position), les onze mots-clés réservés, et tous les exemples (y compris ceux du
dossier \code{exemples/}, tous acceptés). Ils vérifient aussi l'aller-retour
\code{parse(source(p)) == p} sur tous les modèles fournis.

%======================================================================
\section{Question 3 -- Analyse et transformation de l'arbre}

Le travail est découpé en deux passes.

\paragraph{Passe 1 : analyse statique (\code{analyse.py}).} Les instructions sont parcourues
dans l'ordre du texte, avec un seul espace de noms : un nom doit être déclaré avant d'être
utilisé, ce que respectent tous les exemples.
\begin{itemize}
  \item Une \emph{déclaration} est vérifiée, puis ses valeurs connues (tailles, bornes de
        domaine, données des tableaux constants) sont calculées aussitôt et rangées dans la
        table des symboles. Celle-ci connaît les constantes entières, les tableaux constants
        (dimensions, données, type des cases), les types énumérés et leurs valeurs (codes
        $0, 1, \dots$), les variables et les tableaux de variables (dimensions, domaine).
  \item Une \emph{contrainte} est typée une seule fois, sur son texte, avec les indices de
        boucle typés par leur domaine. Le typage renvoie le type (entier, booléen, type
        énuméré), un indicateur \og valeur connue \fg{} et un indicateur \og tranche \fg.
\end{itemize}
Les erreurs détectées sont : nom non déclaré ou déclaré deux fois, y compris un indice de
boucle ; types incompatibles (couleur contre entier, arithmétique sur des booléens,
\code{<} entre booléens, connecteur sur un entier, contrainte non booléenne) ; valeur inconnue
là où une valeur connue est exigée (\code{x + 1 = y}, indice, borne, taille, constante) ;
mauvais nombre d'indices ; indice du mauvais type ; tranche hors d'une contrainte globale ;
fonction inconnue ; mauvais arguments d'une globale ou de \code{abs} ; domaine vide ; données
d'un tableau constant de mauvaise forme. Le choix d'une passe statique a une conséquence
voulue : une erreur de type est signalée même dans une partie qu'un filtre rendrait inutile
(\code{i < 3 -> a[i] = 1} avec \code{a} booléen). Le modèle est faux quelle que soit la
valeur de \code{i}.

\paragraph{Passe 2 : déroulement et simplification (\code{deroulement.py}).} Chaque
\code{forall} est remplacé par ses instances. Les bornes peuvent dépendre des indices
précédents (\code{j in i..N}). Un évaluateur calcule chaque expression dans l'environnement des
indices de boucle. Il renvoie soit une valeur connue (entier, booléen, valeur énumérée), soit
un terme de la représentation intermédiaire (\code{ir.py}) :
\begin{itemize}
  \item des \emph{cases} \code{Cell("grid", (1, 2))} ;
  \item les connecteurs \code{INot}, \code{IAnd}, \code{IOr}, \code{IIff} (\code{=} entre
        booléens) ;
  \item des comparaisons normalisées \code{Cmp}. Entre une variable et une valeur, les
        opérateurs sont \code{= != <= >=} ; entre deux variables, \code{= != < <=} ;
  \item les globales \code{Glob}.
\end{itemize}
Les simplifications appliquées :
\begin{itemize}
  \item valeurs connues calculées : arithmétique, \code{abs}, accès aux tableaux constants,
        comparaisons entre valeurs connues ;
  \item évaluation paresseuse de gauche à droite de \code{and}, \code{or}, \code{->} : si
        l'opérande gauche connu suffit, le droit n'est pas évalué ;
  \item connecteurs simplifiés quand un opérande est connu (\code{faux and e} donne faux,
        etc.) ;
  \item négation poussée jusqu'aux feuilles : De Morgan, \code{not (x = 2)} donne \code{x != 2},
        \code{not (x < y)} donne \code{y <= x} ;
  \item \code{x > y} réécrit en \code{y < x}, \code{x < c} en \code{x <= c-1}, \code{c < x}
        en \code{x >= c+1} ;
  \item tranches remplacées par la liste de leurs cases, ligne par ligne ;
  \item arguments connus des globales simplifiés (un argument vrai dans \code{atMostOne}
        force les autres à faux, deux valeurs égales dans \code{allDifferent} donnent faux,
        une valeur $c$ dans \code{allDifferent} donne \code{x != c} pour chaque case $x$).
\end{itemize}
Le court-circuit fait des conditions connues des filtres : dans
\code{i < N -> grid[i + 1][1] != grid[i][1]}, l'accès \code{grid[N + 1][1]} n'est jamais
évalué. Un accès hors bornes effectivement évalué est une erreur, détectée dans cette passe
puisqu'elle dépend des valeurs, avec les valeurs des indices de boucle dans le message :
\begin{center}
\code{l'indice 1 de 'a' vaut 4, hors de 1..3 (dans 'a[(i + 1)]') (pour i = 3)}
\end{center}
Une contrainte qui se simplifie en vrai disparaît. Une contrainte qui se simplifie en faux
donne un avertissement (\og contrainte toujours fausse \fg) et une clause vide dans la FNC.

\paragraph{Tests (\code{tests/test\_analyse.py}).}
\begin{itemize}
  \item La collection \code{exemples/erreurs/} compte 33 modèles erronés ; chacun annonce en
        première ligne le message attendu. Chaque modèle doit être rejeté par une
        \code{CompileError} contenant ce message, jamais par une autre exception Python. En
        ligne de commande, le code de retour est 1 et il n'y a pas de trace Python.
  \item De petits modèles dont on connaît le résultat déroulé et simplifié (normalisations,
        valeurs connues, tranches, globales).
  \item Des accès hors bornes protégés par une condition fausse, qui ne doivent provoquer
        aucune erreur, et d'autres non protégés ou protégés trop tard (\code{a[i+1] and i < 3}),
        qui doivent en provoquer une.
\end{itemize}

%======================================================================
\section{Question 4 -- Traduction en FNC}

\paragraph{Encodage direct.} Une variable de domaine $\{v_1,\dots,v_k\}$ devient $k$
variables booléennes $x_{v}$, avec la clause \og au moins une \fg{} et les $k(k-1)/2$ clauses
\og au plus une \fg. Pour \code{var x : 1..3;}, on obtient exactement les 4 clauses de
l'énoncé (test). Les valeurs énumérées sont numérotées dans l'ordre de déclaration.

\paragraph{Négations.} La question 3 a déjà poussé la négation jusqu'aux comparaisons
(\code{not (x = y)} est devenu \code{x != y}). Il ne reste que la négation d'un booléen, d'une
égalité entre booléens ou d'une contrainte globale : on traduit l'expression, puis on nie la
formule obtenue par De Morgan.

\paragraph{Comparaisons.} Toutes les traductions s'appuient sur les clauses \og exactement
une \fg{} des domaines. Notons $k = |D(x)|$.
\begin{center}
\begin{tabular}{lll}
\toprule
comparaison & traduction retenue & taille \\
\midrule
\code{x = c}, \code{x != c} & $x_c$, $\neg x_c$ (faux, vrai si $c \notin D(x)$) & 1 littéral \\
\code{x <= c} & $\bigvee_{v \le c} x_v$ & 1 clause \\
\code{x >= c} & $\bigvee_{v \ge c} x_v$ & 1 clause \\
\code{x = y} & $\bigwedge_{v} (\neg x_v \vee y_v)$ & $k$ clauses binaires \\
\code{x != y} & $\bigwedge_{v} (\neg x_v \vee \neg y_v)$ & $k$ clauses binaires \\
\code{x < y}, \code{x <= y} & $\bigwedge_{v} (\neg x_v \vee \bigvee_{w > v} y_w)$
  (resp. $w \ge v$) & $k$ clauses, longueur $\le k+1$ \\
\bottomrule
\end{tabular}
\end{center}
\begin{itemize}
  \item \code{x <= c} peut s'écrire $\bigvee_{v\le c} x_v$ (une clause) ou
        $\bigwedge_{v > c}\neg x_v$ (des clauses unitaires). Au sommet d'une contrainte, les
        deux conviennent. Sous une disjonction (\code{b or x <= 2}), la première se fond dans
        la clause englobante ($b \vee x_1 \vee x_2$), alors que la seconde demanderait une
        variable auxiliaire. Nous retenons donc la forme disjonctive, qui donne toujours une
        seule clause.
  \item \code{x = y} s'écrirait naturellement $\bigvee_v (x_v \wedge y_v)$, qui demande une
        variable auxiliaire par valeur. La forme \og $x_v \to y_v$ pour tout $v$ \fg{} n'en
        demande aucune, grâce à \og exactement une \fg{} sur $y$. Quand $v \notin D(y)$, la
        clause se réduit à $\neg x_v$.
  \item \code{x < y} s'écrirait $\bigvee_{v<w}(x_v\wedge y_w)$, soit $O(k^2)$ conjonctions
        et autant de variables auxiliaires. Nous préférons $k$ clauses \og si $x$ vaut $v$,
        alors $y$ prend une valeur $> v$ \fg. Les clauses sont plus longues (jusqu'à $k+1$
        littéraux), mais il n'y a aucune variable auxiliaire et leur nombre est linéaire.
\end{itemize}
Les tests \code{test\_comparaisons} vérifient ces tailles, et la validation aléatoire de la
question 1, étendue aux domaines finis (3 entiers de domaines \code{1..3}, \code{1..3},
\code{2..4}, deux énumérés, constantes hors domaine, toutes les comparaisons, \code{=} entre
booléens, \code{allDifferent} avec constantes), vérifie leur correction.

\paragraph{Contraintes globales.} \code{atLeastOne} est la disjonction des arguments.
\code{atMostOne} utilise l'encodage deux à deux ($n(n-1)/2$ clauses, aucune variable
auxiliaire), et \code{exactlyOne} la conjonction des deux. \code{allDifferent} est traduite,
pour chaque valeur $v$, par \code{atMostOne} des $x_v$. Ces traductions sont des formules
propositionnelles comme les autres, ce qui permet de les utiliser sous une négation : pour
\code{not atMostOne(a, b, c)}, la négation des clauses donne $\bigvee_{i<j} (l_i \wedge l_j)$,
avec une variable auxiliaire par paire (test).

\paragraph{Non traité.} Nous n'avons pas traité le compteur séquentiel de Sinz (\og pour
aller plus loin \fg). L'encodage deux à deux suffit aux tailles rencontrées.

%======================================================================
\section{Question 5 -- Jonction avec la partie I}

\code{solveur.py} écrit la FNC dans un fichier DIMACS : celui de \code{--dimacs}, ou un
fichier temporaire. Ce fichier est relu par \code{lire\_dimacs} de la partie I et résolu par
\code{dpll\_h} avec l'heuristique H4 dynamique, retenue à la question 5 de la partie I. Le
fichier DIMACS est donc la seule interface entre le compilateur et le solveur. Le modèle
renvoyé est complété (\code{completer}) et vérifié sur la FNC (\code{verifier\_modele}). La
valeur de chaque variable est ensuite retrouvée : pour une variable à domaine fini, c'est
l'unique $v$ tel que $x_v$ est vraie. La solution décodée est enfin vérifiée sur les
contraintes déroulées par un évaluateur de l'IR (\code{ir.evaluer}), indépendant de la FNC.

\begin{center}
\begin{tabular}{lrrrrl}
\toprule
modèle & cases & contraintes déroulées & variables & clauses & réponse \\
\midrule
\code{pigeonhole.csp} & 4 & 1 & 12 & 34 & insatisfiable \\
\code{queens.csp} & 8 & 281 & 64 & 736 & satisfiable \\
\code{coloring.csp} & 10 & 17 & 30 & 87 & satisfiable \\
\code{sudoku.csp} & 81 & 57 & 729 & 10\,317 & satisfiable \\
\bottomrule
\end{tabular}
\end{center}
Les réponses sont celles du tableau 1 de l'énoncé. Aucun de ces modèles n'introduit de
variable auxiliaire. Le Sudoku est résolu en 2~s environ, les autres en moins de 0,05~s.

\paragraph{Vérification des solutions (\code{tests/test\_exemples.py}).} Les vérificateurs
lisent uniquement la sortie texte de \code{python csp.py m.csp --solve} :
\begin{itemize}
  \item huit reines : colonnes deux à deux distinctes, pas de diagonale commune ;
  \item coloration : les 15 arêtes du graphe de Petersen relient des couleurs différentes, et
        \code{color[1] = Red}, \code{color[2] = Green} ;
  \item Sudoku : lignes, colonnes et carrés sont des permutations de 1..9, et les cases
        données sont respectées. Ces données sont relues dans le fichier du modèle.
\end{itemize}

%======================================================================
\section{Question 6 -- Modéliser : le modèle compte plus que le solveur}

Les programmes sont dans \code{q6\_meowdoku/}, et les règles dans
\code{q6\_meowdoku/regles/}. Les grilles $N \times N$ à $N$ régions sont produites par le
générateur de la partie I, avec trois graines par taille.

\paragraph{6.1 Conversion.} \code{grille\_vers\_csp.py} lit une grille au format de la partie I
(\code{lire\_grille}) et écrit les déclarations de l'énoncé : \code{N}, \code{M}, \code{K},
\code{enum Zone} et \code{const region[N][M]}. Avec \code{--regles}, il recopie le fichier de
règles à la suite. Un nom de région qui n'est pas un identifiant, ou qui est déjà pris par un
mot-clé ou par les règles (une région nommée \code{K} ou \code{N} apparaît dès 11 régions),
est préfixé par \code{Z\_}.

\paragraph{6.2 Au plus un chat par région.} Pour deux cases distinctes d'une même région,
prises dans l'ordre ligne par ligne, on interdit qu'elles portent toutes deux un chat :
\begin{verbatim}
forall i1 in 1..N, j1 in 1..M, i2 in 1..N, j2 in 1..M {
    constraint (i1 < i2 or (i1 = i2 and j1 < j2)) and region[i1][j1] = region[i2][j2]
               -> not (cat[i1][j1] and cat[i2][j2]);
}
\end{verbatim}
La condition ne porte que sur des valeurs connues : elle filtre les paires à la compilation,
et chaque instance retenue donne une clause binaire. \code{verifier.py} est un vérificateur
indépendant : il ne lit que la grille et la sortie texte de \code{csp.py --solve}, retrouve les
chats (\code{row[z]}, \code{col[z]} ou \code{cat[i][j] = true}), contrôle pour le premier modèle
que le chat de $z$ est bien dans $z$, puis applique le vérificateur de la partie I. Le script
\code{resoudre\_grilles.py} résout la grille de l'énoncé et 15 grilles de tailles 4 à 8 avec
les quatre fichiers de règles : les 64 solutions sont validées. Sur la grille de l'énoncé, les
deux modèles redonnent la solution de la partie I.

\paragraph{Protocole des mesures (6.3, 6.4).} \code{experiences.py} compile chaque modèle,
écrit le DIMACS et le résout avec l'instrumentation de la partie I (\code{mesures.mesurer} :
décisions, retours arrière, propagations, temps CPU), avec une limite de 60~s. La campagne
compte 156 mesures, sur 3 grilles par taille. Les tableaux donnent la médiane sur les trois
grilles ; « TO » compte les dépassements de la limite. Toutes les mesures terminées sont
SAT, avec un modèle valide.

\paragraph{6.3 Comparaison des deux modèles.}
\begin{center}
\begin{tabular}{r rrrr rrrr}
\toprule
 & \multicolumn{4}{c}{premier modèle (\code{row}, \code{col})} & \multicolumn{4}{c}{second modèle (\code{cat})} \\
\cmidrule(lr){2-5}\cmidrule(lr){6-9}
$N$ & var. & clauses & décisions & CPU (s) & var. & clauses & décisions & CPU (s) \\
\midrule
6  &  72 &  2\,052 &  48 & 0,02 &  36 &    267 &  23 & 0,00 \\
8  & 128 &  6\,848 &  96 & 0,11 &  64 &    642 &  47 & 0,01 \\
10 & 200 & 17\,300 & 162 & 0,50 & 100 & 1\,268 &  78 & 0,02 \\
12 & 288 & 36\,720 & 241 & 1,88 & 144 & 2\,215 & 119 & 0,06 \\
14 & 392 & 69\,188 & 334 & 8,43 & 196 & 3\,597 & 165 & 0,18 \\
\bottomrule
\end{tabular}

\smallskip
{\small Heuristique H4 dynamique (celle de la partie I). Compilation à $N=14$ : 7 à 9~s
pour le premier modèle, 0,4~s pour le second.}
\end{center}
Le second modèle est plus rapide, et l'écart grandit avec $N$ : un facteur 11 à $N = 8$,
47 à $N = 14$. Le compilateur et
le solveur étant les mêmes, l'écart vient de la formule produite.
\begin{itemize}
  \item \textbf{Taille.} Dans le premier modèle, la règle de voisinage porte sur tous les
        couples de régions et toutes les positions : environ $\frac{K^2}{2}\cdot N M \cdot 4$
        instances, soit $O(N^4)$ clauses de 4 littéraux (69\,188 clauses à $N = 14$, 19 fois
        plus que le second modèle). La règle « le chat de $z$ est dans $z$ » ajoute
        $K(NM - |z|)$ clauses binaires. Le second modèle n'a que des clauses binaires en
        $O(N^3)$ et $2N$ clauses « au moins un ». Or chaque appel de notre DPLL recopie la
        formule : le coût d'un nœud est proportionnel au nombre de clauses.
  \item \textbf{Propagation.} Le premier modèle fait aussi deux fois plus de décisions. Une
        position de chat y est une conjonction $\code{row}[z] = i \wedge \code{col}[z] = j$.
        Fixer la ligne d'un chat ne dit rien de ses voisins tant que sa colonne est inconnue,
        car une clause de voisinage de 4 littéraux ne propage qu'une fois trois d'entre eux
        faux. Dans le second modèle, poser un chat sur $(i,j)$ élimine aussitôt, par des
        clauses binaires, toute sa ligne, sa colonne, sa région et ses voisins diagonaux. Les
        clauses « exactement un par ligne et par colonne » forcent en outre la dernière case
        libre d'une ligne. Le premier modèle n'a pas d'équivalent : \code{allDifferent} ne
        donne que des clauses « au plus un » par valeur.
  \item Avec l'heuristique aléatoire H0, le premier modèle dépasse la limite sur 3 des 6
        grilles de tailles 12 et 14, alors que le second les résout toutes (médiane 0,14~s à
        $N = 14$, au plus 9,4~s).
\end{itemize}

\paragraph{6.4 Contrainte redondante : au moins un chat par région.} Aucune tranche ne
désigne les cases d'une région. Nous l'écrivons dans le langage, sans l'étendre, avec un
tableau auxiliaire :
\begin{verbatim}
array inzone[Zone][N][M] : bool;
forall z in Zone, i in 1..N, j in 1..M {
    constraint inzone[z][i][j] = (region[i][j] = z and cat[i][j]);
}
forall z in Zone { constraint atLeastOne(inzone[z][*][*]); }
\end{verbatim}
À la compilation, la condition connue \code{region[i][j] = z} fait de chaque contrainte
\code{not inzone[z][i][j]} (case hors de $z$) ou \code{inzone[z][i][j] = cat[i][j]} (case de
$z$). La tranche \code{inzone[z][*][*]} désigne alors exactement les cases de $z$ : après
propagation, la clause « au moins un » est celle de la partie I. Le prix est de $K N M$
variables auxiliaires, dont $K N M - N M$ fixées à faux par des clauses unitaires.
\begin{center}
\begin{tabular}{r rrr rrr rr}
\toprule
 & \multicolumn{3}{c}{sans la règle} & \multicolumn{3}{c}{avec la règle} & \multicolumn{2}{c}{CPU H4 (s)} \\
\cmidrule(lr){2-4}\cmidrule(lr){5-7}\cmidrule(lr){8-9}
$N$ & var. & déc. H1 & CPU H1 (s) & var. & déc. H1 & CPU H1 (s) & sans & avec \\
\midrule
6  &  36 &       6 &          0,00 &    252 &  4 &       0,03 & 0,00 & 0,03 \\
8  &  64 &       6 &          0,00 &    576 &  6 &       0,21 & 0,01 & 0,19 \\
10 & 100 &      78 &          0,03 & 1\,100 & 15 &       0,94 & 0,02 & 0,91 \\
12 & 144 & 4\,102 &          1,69 & 1\,872 & 14 &       3,67 & 0,06 & 3,07 \\
14 & 196 &     868 &          0,35 & 2\,940 & 16 &      14,71 & 0,18 & 13,75 \\
16 & 256 & 38\,595 & 53 (2 TO) & 4\,352 & -- & (3 TO) & 0,33 & (2 TO) \\
18 & 324 & 61\,313 & 49 (2 TO) & 6\,156 & -- & (3 TO) & 0,62 & (3 TO) \\
20 & 400 & 69\,710 & 45 (3 TO) & 8\,400 & -- & (3 TO) & 1,23 & (3 TO) \\
\bottomrule
\end{tabular}
\end{center}
Deux effets s'opposent.
\begin{itemize}
  \item \textbf{La redondance élague l'arbre.} Avec H1, qui pose un chat sur la première case
        libre et doit souvent revenir en arrière, la règle fait passer le nombre de décisions
        de 4\,102 à 14 à $N = 12$, et les retours arrière de 4\,094 à 6. Sans elle, une région
        dont toutes les cases sauf une sont exclues n'oblige à rien ; avec elle, la dernière
        case est forcée par propagation unitaire, et une région vidée provoque un conflit
        immédiat au lieu d'être découverte bien plus bas dans l'arbre. Avec H4, qui ne
        revient déjà jamais en arrière sur ces grilles, la règle ne change pas le nombre de
        décisions.
  \item \textbf{Les variables auxiliaires coûtent cher à notre solveur.} Les $K N M$
        variables \code{inzone} sont presque toutes affectées par propagation unitaire, et
        chaque propagation recopie la formule. Le temps croît donc comme
        $K N M \times |F|$ : 13,7~s à $N = 14$ contre 0,18~s, quelle que soit l'heuristique.
        À partir de $N = 16$, la plupart des résolutions dépassent la limite de 60~s,
        certaines avant même la première décision.
\end{itemize}
Dans notre implémentation, la règle redondante ne paie donc que lorsque la recherche est
mauvaise (H1). Ce coût vient de notre écriture, qui passe par des auxiliaires, et du coût
linéaire d'une propagation dans notre DPLL ; ce n'est pas un défaut de la redondance
elle-même. Un solveur à littéraux surveillés rendrait ces propagations presque gratuites, et
une liste en compréhension (question optionnelle, non traitée) donnerait directement la
clause, sans auxiliaire.

\paragraph{6.5 Grilles quelconques.} Pour $N$, $M$ et $K$ libres, \code{modele2\_general.csp}
reprend les règles de la partie I :
\begin{itemize}
  \item \emph{au plus} un chat par ligne et par colonne (\code{atMostOne} : « exactement
        un » n'est plus vrai si $K < N$) ;
  \item pas deux chats voisins en diagonale (les voisins sur une même ligne ou une même
        colonne sont déjà exclus) ;
  \item exactement un chat par région : la règle 6.2, plus la règle 6.4, qui n'est plus
        redondante.
\end{itemize}
\code{exhaustif.py} engendre des grilles aléatoires ($N, M \in [2,5]$, $1 \le K \le 6$ régions
connexes, sans garantie de solution). Il compare la réponse du compilateur à une recherche
exhaustive indépendante : il essaie tous les choix d'une case par région et garde ceux que le
vérificateur de la partie I accepte. Sur 300 grilles (141 SAT, 159 UNSAT), les 300 réponses
coïncident, et les 141 solutions sont validées par \code{verifier.py}.

\paragraph{6.6 Conclusion.} Les deux leviers comptent, mais pas au même niveau.
\begin{itemize}
  \item \textbf{Le choix du modèle fixe l'ordre de grandeur.} À $N = 14$, la \emph{meilleure}
        heuristique sur le premier modèle (H4 : 8,4~s, 7 à 9~s de compilation en plus) est
        60 fois plus lente, en médiane, que la \emph{pire} sur le second (H0 : 0,14~s). Le premier modèle
        est plus gros, en $O(N^4)$ clauses contre $O(N^3)$, et propage moins. Aucune
        heuristique ne rattrape cela, puisque le coût de chaque nœud est proportionnel à la
        taille de la formule.
  \item \textbf{L'heuristique compte ensuite, et peut être catastrophique.} Sur le second
        modèle, H1 dépasse la limite sur 7 des 9 grilles de tailles 16 à 20, que H4 résout en
        moins de 1,3~s. En partie I, les heuristiques apportaient un facteur 4 à 37.
  \item \textbf{Les deux interagissent.} H1, la meilleure sur le premier modèle, est la pire
        sur le second. Une contrainte redondante, qui est un choix de modélisation, divise par
        300 le nombre de décisions de H1, mais elle est sans effet pour H4.
\end{itemize}
Dans nos expériences, le choix du modèle est donc ce qui a le plus d'effet : il détermine la
taille de la formule et la force de la propagation unitaire, c'est-à-dire le terrain sur
lequel l'heuristique travaille. Les écarts entre heuristiques sont importants, mais ils se
jouent à l'intérieur d'un modèle donné.

%======================================================================
\section*{Ce qui n'a pas été traité}
Nous n'avons pas traité les parties \og pour aller plus loin \fg{} (positions ligne et colonne
dans les messages d'erreur, compteur séquentiel de Sinz) ni les questions optionnelles.
Les messages d'erreur sémantique ne donnent donc pas de position, mais ils citent
l'expression en cause et, au déroulement, les valeurs des indices de boucle.

%======================================================================
\section*{Usage des assistants IA}
 Nous avons d'abord défini notre approche d'implémentation.
 Face aux premiers bugs rencontrés, nous nous sommes appuyés sur Claude Code pour les corriger. 
 Nous avons ensuite conçu les tests en collaboration avec l'outil, et vérifier nous-mêmes les résultats obtenus. 
 Enfin, Claude Code a généré le rapport, que nous avons intégralement relu et corrigé.

\end{document}
